\color{unimain}\section[Сервисные функции]{Сервисные функции}\label{sec:servfuncs}
\color{black}



%%%%%%%%%%%%%%%%%%%%%%%%%%%%%%%%%%%%%%%%%%%%%% ЖУРНАЛ СОБЫТИЙ %%%%%%%%%%%%%%%%%%%%%%%%%%%%%%%%%%%%%%%%%%%%%%%%%%
\needspace{4\baselineskip}
\color{unimain}{\subsection{Журнал событий}}
\color{black}

\begin{enumerate}

\item
устройстве реализована функция записи сигналов дискретных входов, работы защит, автоматики и сервисных сигналов в журнал событий.
\item
Максимальное число записей журнала - 20000. Каждая запись в журнал событий содержит в себе следующую информацию:
\begin{itemize}
    \item дату и время возникновения события;
    \item идентификатор события;
    \item состояние дискретного или значение аналогового сигнала в момент события;
    \item качество зарегистрированного параметра;
    \item качество времени;
\end{itemize}
\item
Данные о событиях сохраняются в энергонезависимой памяти. При заполнении журнала событий перезапись осуществляется от более старых событий к новым, то есть циклически. Удаление информации из журнала событий пользователем невозможно.

\end{enumerate}

%%%%%%%%%%%%%%%%%%%%%%%%%%%%%%%%%%%%%%%%%%%%%% РЕЖИМ ТЕСТИРОВАНИЯ %%%%%%%%%%%%%%%%%%%%%%%%%%%%%%%%%%%%%%%%%%%%%%%%%%
\needspace{4\baselineskip}
\color{unimain}{\subsection{Режим тестирования}}
\color{black}

\begin{enumerate}

\item
Режим тестирования предусмотрен для возможности оперативной проверки элементов/функций, таких как: ФСУ, измерительные органы, дискретные входы и выходные реле устройства.

\end{enumerate}

%%%%%%%%%%%%%%%%%%%%%%%%%%%%%%%%%%%%%%%%%%%%%% Синхронизация времени %%%%%%%%%%%%%%%%%%%%%%%%%%%%%%%%%%%%%%%%%%%%%%%%%%
\needspace{4\baselineskip}
\color{unimain}{\subsection{Синхронизация времени}}
\color{black}

\begin{enumerate}

\item
Аппаратные средства устройства предусматривает возможность синхронизации внутренних часов реального времени ЦП устройства одним или несколькими из способов: в соответствии с протоколами
NTP(SNTP), PTP, секундным импульсом 1PPS, средствами стандартизированного протокола передачи данных.

\end{enumerate}

%%%%%%%%%%%%%%%%%%%%%%%%%%%%%%%%%%%%%%%%%%%%%% Система самодиагностики %%%%%%%%%%%%%%%%%%%%%%%%%%%%%%%%%%%%%%%%%%%%%%%%%%
\needspace{4\baselineskip}
\color{unimain}{\subsection{Система самодиагностики}}
\color{black}

\begin{enumerate}

\item
Система самодиагностики устройства выполняет тесты в полном объеме при включении устройства, а также непрерывно в фоновом режиме.
Системой самодиагностики устройства в автоматическом режиме контролируются следующие параметры:
\begin{itemize}
\item состояние аппаратной части устройства, в том числе: АЦП модулей ввода аналоговых сигналов, блока питания, ОЗУ, ПЗУ, процессора, модулей дискретных входов, модулей дискретного управления, вспомогательных модулей;
\item температурный режим устройства;
\item состояние синхронизации времени;
\item сохранность исполняемого программного кода (целостность ПО);
\item состояние измерительных цепей;
\item исправность используемых каналов связи.
\end{itemize}

\item
Система самодиагностики фиксирует критические неисправности («Неисправность») и некритические ошибки («Ошибка»). По факту выхода из режимов критической неисправности и предупредительной сигнализации производится запись в журнале событий. Подробное описание системы самодиагностики приводится в документации ЮТКБ.656122.609 РЭ «Устройство защиты и автоматики серии «ЮНИТ-М300». Руководство по эксплуатации общее».

\item
В случае появления критической неисправности активируется внутренний сигнал «Критическая неисправность», производится запись в журнал событий, замыкаются контакты реле IRF, прекращается работа алгоритмов основных функций в составе устройства, это состояние устройства сигнализируется свечением светодиода «Неисправность» и отсутствием свечения светодиода «Готовность». Для выхода из режима критической неисправности требуется перезапуск устройства.

\item
В случае появления некритической ошибки активируется внутренний сигнал «Некритическая ошибка», производится запись в журнал событий, алгоритмы основных функций в составе устройства продолжают выполняться, это состояние устройства сигнализируется свечением светодиода «Неисправность» и отсутствием свечения светодиода «Готовность». Устройство может выйти из режима предупредительной сигнализации без активных действий со стороны пользователя по факту пропадания критериев режима.

\end{enumerate}

%%%%%%%%%%%%%%%%%%%%%%%%%%%%%%%%%%%%%%%%%%%%%% Аутентификация и авторизация пользователей %%%%%%%%%%%%%%%%%%%%%%%%%%%%%%%%%%%%%%%%%%%%%%%%%%
\needspace{4\baselineskip}
\color{unimain}{\subsection{Аутентификация и авторизация пользователей}}
\color{black}

\begin{enumerate}

\item
Для возможности работы с устройством пользователь должен пройти аутентификацию на устройстве. Аутентификация проводится по имени и паролю, указанными пользователем. 
\item
Функции безопасности устройства предъявляют следующие требования к паролям пользователей:
\begin{itemize}
    \item пароль может содержать буквы и/или цифры;
    \item длина пароля не может быть меньше семи символов. 
\end{itemize}
\item
Хранение паролей осуществляется в памяти устройства в нечитаемом виде, не позволяющем восстановить пароль.
\item
Функции безопасности не дают возможность повторного применения ранее использованных паролей. По умолчанию глубина хранения ранее заданных пользователем паролей в устройстве составляет четыре пароля.
\item
В процессе ввода пароля пользователю отображается только вводимый символ, ранее введенные символы не отображаются в открытом виде.
\item
Функции безопасности устройства ограничивают количество неуспешных попыток аутентификации за установленный период времени. Количество попыток можно задать в диапазоне от 1 до 9 (по умолчанию – 3)
\item
В случае превышения количества неуспешных попыток аутентификации:
\begin{itemize}
    \item выполняется блокировка доступа пользователя к устройству. Длительность блокировки определяется в настройках устройства пользователем с ролью Администратор в диапазоне от 1 минуты до 9999 минут (по умолчанию – 30 минут);
    \item формируется соответствующее событие журнала событий безопасности;
    \item формируется сигнализация для АСУ ТП о превышении заданного количества неуспешных попыток доступа. 
\end{itemize}
\item
В случае, когда количество неуспешных попыток аутентификации пользователя больше нуля, но меньше максимально допустимого, оно будет сброшено для данного пользователя по истечении времени сброса неуспешных попыток (5 минут).
\item
Блокировка возможности аутентификации пользователя снимается в случае:и:
\begin{itemize}
    \item истечения времени блокировки входа;
    \item принудительной разблокировки пользователем с ролью Администратор. 
\end{itemize}

\item
Пользователь может быть заблокирован принудительно пользователем с ролью Администратор. Разблокирование данного пользователя может осуществить так же только пользователь с ролью Администратор.
\item
Все настраиваемые параметры безопасности (время неактивности пользователя, количество безуспешных попыток входа, время блокировки входа) разрешены для редактирования пользователям с ролью офицера безопасности. При сбросе устройства к заводским настройкам, параметры безопасности сохраняются.

%\noindent
%\colorbox{yellow!40}{
 %\begin{minipage}{\dimexpr\linewidth-2\fboxsep}
    %{\faExclamationTriangle} При утере пароля учетной записи пользователя с ролью Администратор, восстановление административных прав возможно только путем перепрограммирования устройства специалистом предприятия-изготовителя, что повлечет за собой сброс настроечных параметров на заводские.
  %\end{minipage}
%}

% ПРЕДУПРЕЖДЕНИЕ - Несоблюдение указанных требований может нанести ущерб оборудованию или нарушению работы энергосистемы
%\noindent
%\fboxrule=1.2pt % задаём толщину рамки (жирная линия)
%\noindent
%\fbox{%
  %\begin{minipage}{\dimexpr\linewidth-2\fboxrule-2\fboxsep}
    %{\hspace*{0.5em}\textcolor{unired}{\Large\faExclamationTriangle}}\hspace{0.5em}
    %При утере пароля учетной записи пользователя с ролью Администратор, восстановление административных прав возможно только путем перепрограммирования устройства специалистом предприятия-изготовителя, что повлечет за собой сброс настроечных параметров на заводские.
  %\end{minipage}%
%}

% ОПАСНОСТЬ - Несоблюдение указанных требований может повлечь угрозу здоровью и жизни
%\noindent
%\fboxrule=1.2pt % задаём толщину рамки (жирная линия)
%\noindent
%\fbox{%
  %\begin{minipage}{\dimexpr\linewidth-2\fboxrule-2\fboxsep}
    %{\hspace*{0.5em}\textcolor{red}{\Large\faBolt}}\hspace{0.5em}
    %При утере пароля учетной записи пользователя с ролью Администратор, восстановление административных прав возможно только путем перепрограммирования устройства специалистом предприятия-изготовителя, что повлечет за собой сброс настроечных параметров на заводские.
  %\end{minipage}%
%}

% ИНФОРМАЦИЯ - 
\noindent
\fboxrule=1.2pt % задаём толщину рамки (жирная линия)
\noindent
\fbox{%
  \begin{minipage}{\dimexpr\linewidth-2\fboxrule-2\fboxsep}
    {\hspace*{0.5em}\textcolor{unimain}{\Large\faInfoCircle}}\hspace{0.5em}
    \hyphenpenalty=10000 При утере пароля учетной записи пользователя с ролью Администратор, восстановление административных прав возможно только путем перепрограммирования устройства специалистом предприятия-изготовителя, что повлечет за собой сброс настроечных параметров на заводские.
  \end{minipage}%
}


\end{enumerate}

%%%%%%%%%%%%%%%%%%%%%%%%%%%%%%%%%%%%%%%%%%%%%% Блок интерфейсный (ЮНИТ-ИЧМ) %%%%%%%%%%%%%%%%%%%%%%%%%%%%%%%%%%%%%%%%%%%%%%%%%%
\needspace{4\baselineskip}
\color{unimain}{\subsection{Блок интерфейсный (ЮНИТ-ИЧМ)}}
\color{black}

\begin{enumerate}

\item
Для местного управления и контроля за состоянием устройства предусматривается подключение блока интерфейсного «ЮНИТ-ИЧМ» производства ООО «Юнител Инжиниринг». Блок подключается к «ЮНИТ-М300» посредством патч-корда с разъемами RJ45. В устройстве «ЮНИТ-М300» для подключения к блоку может использоваться отдельный порт X5, расположенный на модуле центрального процессора или любой из задействованных портов связи Ethernet (X1, Х2, Х3, X4).
\item
Более подробная информация по работе с устройством «ЮНИТ-ИЧМ» приведена в документации RU.37182817.00001.02.34.01 «Программное обеспечение мониторинга и параметрирования «ЮНИТ-СЕРВИС». Руководство пользователя».

\item
Подробное описание блока приведено в документе ЮТКБ.656132.701 РЭ – «Блок интерфейсный «ЮНИТ-ИЧМ». Руководство по эксплуатации». Блок интерфейсный «ЮНИТ-ИЧМ» является опциональным для заказа.

\end{enumerate}